\documentclass{article}
\usepackage{amsmath, amssymb, amsthm}
\usepackage{geometry}
\geometry{margin=1in}

\newtheorem{theorem}{Theorem}
\newtheorem{corollary}{Corollary}

\title{Random Walks -- Lecture Notes}
\author{}
\date{}

\begin{document}
\maketitle

%% -------------------------------------------------------
%% Image 1: Pólya's Urn Scheme
%% -------------------------------------------------------
\section{P\'olya's Urn Scheme}

\textbf{Setup.} Initially the urn contains $r$ red balls and $g$ green balls.
At each step, draw one ball uniformly at random, observe its color, then
return it together with $c$ additional balls of the same color.

\textbf{Notation.}
\begin{align*}
  G_n &= \text{number of green balls after the }n\text{-th draw (and add)},\\
  R_n &= \text{number of red balls after the }n\text{-th draw (and add)}.
\end{align*}

\textbf{Proportion process.} Define
\[
  X_n = \frac{G_n}{G_n + R_n},
\]
the proportion of green balls after $n$ draws.

%% -------------------------------------------------------
%% Image 2: X_n is a martingale + Convergence theorem
%% -------------------------------------------------------
\section{$X_n$ is a Martingale}

\textbf{Claim.} $(X_n)_{n \ge 0}$ is a martingale with respect to the natural
filtration $(\mathcal{F}_n)$.

\begin{proof}
Conditioning on $\mathcal{F}_n$, the $(n+1)$-st draw is green with probability
$G_n/(G_n+R_n)$ and red with probability $R_n/(G_n+R_n)$. Therefore
\begin{align*}
  E[X_{n+1} \mid \mathcal{F}_n]
  &= \frac{G_n}{G_n+R_n} \cdot \frac{G_n+c}{G_n+R_n+c}
   + \frac{R_n}{G_n+R_n} \cdot \frac{G_n}{G_n+R_n+c} \\[4pt]
  &= \frac{G_n(G_n+c) + G_n R_n}{(G_n+R_n)(G_n+R_n+c)} \\[4pt]
  &= \frac{G_n(G_n + c + R_n)}{(G_n+R_n+c)(G_n+R_n)} \\[4pt]
  &= \frac{G_n}{G_n+R_n} = X_n. \qedhere
\end{align*}
\end{proof}

\section{Martingale Convergence}

\begin{theorem}[$L^2$ Martingale Convergence Theorem]
If $(X_n)_{n \ge 0}$ is an $L^2$-submartingale satisfying
\[
  \sup_{n} E[X_n^2] < \infty,
\]
then there exists $X_\infty \in L^2$ such that $X_n \to X_\infty$ almost surely.
\end{theorem}

\begin{corollary}[4.7.12]
If $(X_n)_{n \ge 0}$ is a supermartingale with $X_n \ge 0$ for all $n$, then
\[
  X_n \xrightarrow{a.s.} X_\infty \quad \text{and} \quad E[X_\infty] \le E[X_n] \quad \forall\, n.
\]
\end{corollary}

\begin{proof}
Set $Y_n = -X_n$. Then $(Y_n)$ is a submartingale satisfying $Y_n \le 0$, so
$\sup_n E[Y_n^2] = \sup_n E[X_n^2] < \infty$ (since the $X_n$ are bounded below
by $0$ and $E[X_n]$ is decreasing). By the $L^2$ convergence theorem,
$Y_n \to Y_\infty$ a.s., hence $X_n \to X_\infty \ge 0$ a.s.

For the bound on expectations, apply Fatou's lemma:
\[
  E[X_\infty]
  = E\!\left[\liminf_{n} X_n\right]
  \le \liminf_{n}\, E[X_n]
  \le E[X_m] \quad \forall\, m.
\]
Hence $E[X_\infty] \le E[X_n]$ for all $n$.
\end{proof}

\end{document}
